\minitopic{还原论的休谟式动机}还原论者要求，听者接受证言的权利最终来自非证言来源：对人类诚实倾向的归纳、对当前说话者能力的观察、内容与背景知识的一致。其吸引力在于避免把任何断言自动变成证据。陌生人提出极不寻常主张时，我们确实会要求额外支持。 但严格全局还原面临启动问题。儿童尚未积累说话者记录，却必须通过他人学习语言与世界；成人的大部分背景知识本身也来自证言，用这些知识验证新证言可能形成广泛依赖循环。若还原只要求对证言“一般可靠”有证据，主体如何在不依赖证言的情况下建立如此宏大归纳？ 局部还原较温和：每次接受都需有关于当前来源的正面理由。它仍可能过强。问路、课堂常识和普通通知通常在无警报时直接接受，而不先进行可靠性推理。听者对沟通实践的能力可能体现为默认信任加击败敏感，而非显式证明。

\minitopic{非还原论与先验权利}非还原论主张，理解一个真诚断言且无反证，本身给予可被击败的接受权利。Burge把这种权利联系到理性内容保存：表达者把内容呈现为真，听者有初步资格接受\parencite{burge1993}。该立场不说人通常诚实的频率无关，而是不要求每次权利都从个体归纳产生。 默认权利须有边界。广告、匿名谣言、已知利益冲突和明显专业越界会触发审查。无击败不是完全空白：听者需要正常理解、来源看似具备基本能力，传播情形没有欺骗迹象。把这些条件扩张过多，非还原论会悄悄变回还原论。 非还原论还解释沟通为何具有特殊规范。说话者断言$p$不仅造成听者信念，还承担可被质问的承诺。若他只是猜测，应降低语气；若隐瞒关键信息，即使字面真也可能误导。证言的认识力量部分来自这种责任结构。

\minitopic{保证观点}保证观点强调，听者不是仅根据“此人说$p$”作为一个迹象推断$p$，而是接受说话者为$p$作保。说话者邀请听者依赖自己，并承担解释和纠错责任。这能区分人与温度计：两者都传递信息，只有人能以沟通行为承担保证。 反对者指出，录音、档案和匿名文本仍可产生知识，即使原作者无法回应；非人系统也能成为可靠来源。因此保证或许解释面对面证言的特殊价值，却不是所有证言知识的必要条件。可将证言分成沟通性信任与信息性依赖两种模式。 保证关系还可能受权力扭曲。下属被迫为机构结论背书，实际没有审查权限；平台把自动生成文本显示成个人建议，模糊责任主体。真正保证需要主体有相应认识地位和承担修正的能力，不能只看语气。

\minitopic{传递、生成与汇聚}传递论认为，证言像管道：说话者拥有的知识或确证传给听者。若源头无知识，管道不能创造知识。这适合许多教师—学生案例，却在复杂链条中受挑战。 听者可能比说话者拥有更多背景。他听到零散事实后推出说话者未认识的结论，证言成为生成性材料。多个各自不充分的报告也可汇聚成知识。反过来，说话者知道$p$，听者若有强反证或误解含义，传递会失败。 机构报告常混合传递与生成。数据采集者知道局部结果，统计团队形成总体结论，发言人传递组织结论。没有单一人的知识完整沿链移动，但机制可产生群体和听者知识。模型应追踪各环节权限，而非只问第一位说话者是否知道最终命题。

\minitopic{专家的元专长问题}非专家用什么能力识别专家？若必须先掌握领域内容，就不再是非专家；若只看头衔，又容易受假权威欺骗。元专长不是判断所有一阶结论，而是评价训练、同行记录、方法透明、预测校准和对质疑的响应\parencite{goldman2001}。 共识是重要指标，但须分析形成过程。独立研究汇合比共享同一数据错误更有力；开放竞争后的共识比压制异议后的表面一致更有力；成熟领域的稳定共识与新领域的暂时多数权重不同。少数意见也应按证据和记录评价，不因少数身份自动增加或减少可信度。 利益冲突是风险因素，不是逻辑反驳。资助来源可能增加偏差概率，要求披露、独立复核和更强证据；不能仅从“有利益”推出结论为假。反过来，没有财务利益也不保证能力或公正。

\minitopic{同侪分歧的精细结构}真正同侪需要在相关证据、推理能力、认真程度上大致相当。若对方没看关键数据，不构成同侪；若双方各自拥有私人证据，发现分歧本身提示对方可能有自己缺少的信息。现实评估是局部而非全人格的：优秀医生在统计解释上未必与你同侪。 调和主义认为，发现同侪反对是关于自己可能犯错的高阶证据，应降低信心\parencite{christensen2007}。坚定立场担心全面平分会使证据失去作用：每人正因认为自己的推理更好才持结论。合理回应可不是机械折半，而是重新检查最可能分叉的步骤。 分歧还可能有可预见的系统解释。双方价值取向不同、使用不同损失函数或对模型范围理解不同，未必在同一事实命题上冲突。把表面“专家意见不一”拆成数据、方法、概念和价值分歧，常能恢复局部共识。

\minitopic{认识信任与脆弱性}信任使人把一部分核验交给他人，因而创造脆弱性。若只在已经独立验证一切后才“信任”，信任没有认识作用。成熟信任不是消灭风险，而是依据关系和制度使风险可接受，并保留受骗后的问责渠道。 受压迫群体有时对机构持有基于历史记录的合理不信任。简单要求“相信专家”会忽略机构过去的欺骗、排除和不透明。重建信任需要改变程序、代表性和申诉渠道，而不只是传播更多结论。 不信任也可被操纵。制造少量怀疑、夸大个别错误，能让公众把任何共识视为阴谋。认识韧性要求同时保留可问责信任与对系统性怀疑策略的识别。

\minitopic{综合案例：公共卫生建议}卫生机构发布新建议，依据多项研究和监测数据。普通人无法重做试验。合理评估可检查：机构是否说明证据等级；建议与事实陈述是否区分；不确定性是否更新；专家委员会有无利益披露；不同国家独立机构是否汇合。 若建议后来改变，不自动证明先前不可信。新数据、疾病变化或目标权衡都可导致理性更新。机构应解释变化原因；公众应区分“过去证据下合理但后来修正”与“过去隐瞒已知反证”。可修正性是科学信任的组成，不是失败标志。 个人行动还涉及价值与风险。专家提供事实和预测，不替代个人或公共决策中的价值权衡。承认分工可避免两个极端：因专家不决定价值而否认其事实权威，或因其有事实专长而让其独占政策选择。

\begin{tcolorbox}[breakable,colback=white,colframe=blue!30!black,title={证言与专家评估表},fonttitle=\bfseries]
追踪原始说话者、传播链和独立性；区分默认权利与正面可靠性证据；检查能力、诚实、利益与可问责性；判断知识是传递、生成还是汇聚；识别分歧属于数据、方法、概念或价值；说明何种新证据会改变信任。
\end{tcolorbox}
